%% toc: P3 27 27.1,27.2,27.3,27.4
\section{Proofs for Horn preservation and isomorphic ultrapowers}
\label{sec:p03-w03-p03-017-annex}

The ambient set theory is \(\mathsf{ZFC}\).  Cardinal arithmetic uses
\cref{thm:p06-cardinal-arithmetic-laws}; in particular
\(\nu\cdot\nu'=\max(\nu,\nu')\) when one of the two is infinite and neither
is zero, and \((2^\lambda)^\nu=2^{\lambda\cdot\nu}\).  Monotonicity of
cardinal powers in base and exponent is used as in
\cref{sec:p03-w03-p03-018-annex}.

\subsection{Horn formulas}

\begin{annexitem}{annex:p03-w03-p03-017-atomic}{thm:p03-w03-p03-017-atomic}{The congruence and the atomic condition}
\begin{proof}
\emph{Item 1.}  \([\![a\approx a]\!]=I\in F\), and
\([\![a\approx b]\!]=[\![b\approx a]\!]\).  If
\([\![a\approx b]\!],[\![b\approx c]\!]\in F\), then their intersection is in
\(F\) and is contained in \([\![a\approx c]\!]\), so the latter is in \(F\)
by upward closure.  If \(a_m\mathrel{\theta_F}b_m\) for \(m\le n\), then
\(\bigcap_m[\![a_m\approx b_m]\!]\in F\); at each index in this set the
arguments agree, hence so do the values of \(f\) and the truth values of
\(r\).  So \([\![f(\bar a)\approx f(\bar b)]\!]\in F\), and
\([\![r(\bar a)]\!]\in F\) iff \([\![r(\bar b)]\!]\in F\), because each of
these sets, intersected with \(\bigcap_m[\![a_m\approx b_m]\!]\), gives the
same set, and a set belongs to the filter \(F\) iff its intersection with a
member of \(F\) does.  This is
\cite[V\S2, Lemma~2.2, p.~206]{BurrisSankappanavar2012}.

\emph{Item 2.}  Induction on \(t\).  For a variable \(x_j\), both sides are
\(a_j/F\).  For \(t=f(t_1,\ldots,t_n)\), with \(f\) possibly a constant
(\(n=0\)), the induction hypothesis gives
\(t_m^{\prod A_i/F}(\bar a/F)=b_{t_m}/F\), and the interpretation of \(f\)
gives \(f(b_{t_1},\ldots,b_{t_n})/F\); coordinatewise this function is
\(i\mapsto f^{A_i}(t_1^{A_i}(\bar a(i)),\ldots)=b_t(i)\).

\emph{Item 3.}  For \(\alpha=r(t_1,\ldots,t_n)\), item~2 and the definition
give: \(\alpha\) holds at \(\bar a/F\) iff
\([\![r(b_{t_1},\ldots,b_{t_n})]\!]\in F\), and this set equals
\([\![\alpha(\bar a)]\!]\) because \(b_{t_m}(i)=t_m^{A_i}(\bar a(i))\).  For
an equation \(t\approx s\), it holds at \(\bar a/F\) iff \(b_t/F=b_s/F\)
iff \([\![b_t\approx b_s]\!]\in F\), and
\([\![b_t\approx b_s]\!]=[\![(t\approx s)(\bar a)]\!]\).

\emph{Item 4.}  Since \(I\neq\varnothing\), \(\{I\}\) is a proper filter, and
\([\![a\approx b]\!]\in\{I\}\) iff \(a(i)=b(i)\) for all \(i\), iff \(a=b\).
So \(a\mapsto a/\{I\}\) is a bijection.  It commutes with the operations by
the definition of the quotient operations, and \(r(\bar a/\{I\})\) holds iff
\([\![r(\bar a)]\!]=I\), which is the direct-product clause for \(r\).
\end{proof}
\end{annexitem}

\begin{annexitem}{annex:p03-w03-p03-017-maximal}{thm:p03-w03-p03-017-maximal}{The maximal property}
\begin{proof}
Let \(S=[\![\exists x\,\varphi(x,\bar a)]\!]\).  For \(i\in S\) the set
\(W_i=\{c\in A_i:A_i\models\varphi(c,\bar a(i))\}\) is nonempty; for
\(i\notin S\) put \(W_i=A_i\), nonempty by hypothesis.  By the Axiom of
Choice there is \(a\in\prod_iW_i\subseteq\prod_iA_i\).  Then
\([\![\varphi(a,\bar a)]\!]\supseteq S\) by the choice on \(S\), and
\([\![\varphi(a,\bar a)]\!]\subseteq S\) because a witness at \(i\) makes
the existential true at \(i\).  This is
\cite[V\S2, Lemma~2.6, p.~207]{BurrisSankappanavar2012}.
\end{proof}
\end{annexitem}

\begin{annexitem}{annex:p03-w03-p03-017-horn-preservation}{thm:p03-w03-p03-017-horn-preservation}{Horn formulas are preserved under reduced products}
\begin{proof}
Write \(P=\prod_{i\in I}A_i/F\).  The proof follows
\cite[V\S2, Theorem~2.7, pp.~207--210]{BurrisSankappanavar2012} and is by
induction on the number \(m\) of quantifiers in the prefix, simultaneously
for all tuples \(\bar a\).

\emph{A basic Horn formula.}  Let
\(\beta=\beta_1\vee\cdots\vee\beta_k\) with \([\![\beta(\bar a)]\!]\in F\).
\begin{itemize}
  \item If some disjunct is a negated atom \(\neg\alpha_j\) with
        \([\![\alpha_j(\bar a)]\!]\notin F\), then \(P\not\models\alpha_j(\bar
        a/F)\) by item~3 of \cref{thm:p03-w03-p03-017-atomic}, so
        \(P\models\neg\alpha_j(\bar a/F)\) and \(P\models\beta(\bar a/F)\).
  \item Otherwise every negated disjunct \(\neg\alpha_j\) has
        \([\![\alpha_j(\bar a)]\!]\in F\), and so
        \(T=[\![\beta(\bar a)]\!]\cap\bigcap_j[\![\alpha_j(\bar a)]\!]\in F\),
        the intersection being finite.  If \(\beta\) had no atomic disjunct,
        then at every index of \(T\) every disjunct would be false while
        \(\beta\) is true, so \(T=\varnothing\), contradicting
        \(\varnothing\notin F\).  Hence \(\beta\) has exactly one atomic
        disjunct \(\alpha\).  At each index of \(T\) all negated disjuncts
        are false, so \(\alpha\) is true: \(T\subseteq[\![\alpha(\bar a)]\!]\).
        Thus \([\![\alpha(\bar a)]\!]\in F\), and \(P\models\alpha(\bar a/F)\)
        by item~3 of \cref{thm:p03-w03-p03-017-atomic}; so
        \(P\models\beta(\bar a/F)\).
\end{itemize}
This case covers equations in either position, and the false-head case is
the one where the second alternative is impossible.

\emph{A conjunction.}  If \([\![\varphi_1\wedge\cdots\wedge\varphi_r(\bar
a)]\!]\in F\), each \([\![\varphi_s(\bar a)]\!]\) contains this set and lies
in \(F\); by the previous case \(P\models\varphi_s(\bar a/F)\) for each
\(s\).  This settles \(m=0\).

\emph{Quantifier step.}  Let \(\varphi=Q_1y_1\,\psi(y_1,\bar x)\), where
\(\psi\) is a Horn formula with \(m-1\) prefix quantifiers.  If \(Q_1=\forall\)
and \(a\in\prod_iA_i\), then
\([\![\varphi(\bar a)]\!]\subseteq[\![\psi(a,\bar a)]\!]\), so the latter is
in \(F\) and the induction hypothesis gives \(P\models\psi(a/F,\bar a/F)\).
Every element of \(P\) is some \(a/F\), so \(P\models\varphi(\bar a/F)\).  If
\(Q_1=\exists\), \cref{thm:p03-w03-p03-017-maximal} (this is where Choice is
used) gives \(a\) with \([\![\psi(a,\bar a)]\!]=[\![\varphi(\bar a)]\!]\in F\);
the induction hypothesis gives \(P\models\psi(a/F,\bar a/F)\), hence
\(P\models\varphi(\bar a/F)\).
\end{proof}
\end{annexitem}

\begin{annexitem}{annex:p03-w03-p03-017-horn-products}{thm:p03-w03-p03-017-horn-products}{Horn sentences in direct products, reduced products and ultraproducts}
\begin{proof}
\emph{Item 1.}  For a sentence, \([\![\sigma]\!]=\{i:A_i\models\sigma\}\);
apply \cref{thm:p03-w03-p03-017-horn-preservation} with \(n=0\).  If
\(\sigma\) holds in every factor then \([\![\sigma]\!]=I\), which lies in
every filter.  An ultrafilter is a proper filter, so ultraproducts are
included.

\emph{Item 2.}  As \(I\neq\varnothing\), \(\{I\}\) is proper and
\([\![\sigma]\!]=I\in\{I\}\), so \(\prod_iA_i/\{I\}\models\sigma\) by item~1.
By item~4 of \cref{thm:p03-w03-p03-017-atomic} this structure is isomorphic
to \(\prod_iA_i\), and isomorphic structures satisfy the same sentences
(\cref{lem:iso-invariance}).

\emph{Item 3.}  If \(\sigma'\) is logically equivalent to the Horn sentence
\(\sigma\), then \(\sigma\) and \(\sigma'\) have the same truth value in every
structure, so \([\![\sigma']\!]=[\![\sigma]\!]\) and both have the same truth
value in each reduced product.
\end{proof}
\end{annexitem}

\subsection{The construction in Shelah's paper}

Throughout, \(\lambda\) is an infinite cardinal,
\(\mu=\min\{\mu:\lambda^\mu>\lambda\}\), and \(\theta=2^\lambda\).  By
\cref{thm:p03-w03-p03-018-mu}, \(\aleph_0\le\mu\le\lambda\) and
\(\lambda^\nu=\lambda\) for \(1\le\nu<\mu\).  Since \(\mu\) is an infinite
cardinal, \(\nu_1\cdot\nu_2<\mu\) and \(\nu_1+\nu_2<\mu\) whenever
\(\nu_1,\nu_2<\mu\): both are finite when \(\nu_1\), \(\nu_2\) are finite,
and otherwise both are at most \(\max(\nu_1,\nu_2)\).

\begin{annexitem}{annex:p03-w03-p03-017-lemma1}{thm:p03-w03-p03-017-lemma1}{Lemma~1: an independent family}
\begin{proof}
Following \cite[p.~226]{Shelah1971Ultrapowers}, let \(H\) be the set of
pairs \((A,h)\) with \(A\subseteq\lambda\), \(\abs A<\mu\), and \(h\) a
function from a family \(S\) of fewer than \(\mu\) subsets of \(A\) into
\(\mu\).

\emph{Counting.}  For \(1\le\chi<\mu\) there are at most
\(\lambda^\chi=\lambda\) subsets of \(\lambda\) of size \(\chi\) (each is the
range of a function from \(\chi\) to \(\lambda\)), and one of size \(0\);
summing over the at most \(\mu\le\lambda\) cardinals \(\chi<\mu\) gives at
most \(\lambda\) sets \(A\).  For such \(A\), \(\abs{\mathcal P(A)}=
2^{\abs A}\le\lambda^{\abs A}\le\lambda\) (and \(=1\) if \(A=\varnothing\)),
so the families \(S\subseteq\mathcal P(A)\) with \(\abs S=\chi<\mu\) number at
most \(\lambda^\chi\le\lambda\), and at most \(\lambda\) in all.  For such
\(S\), the functions \(S\to\mu\) number \(\mu^{\abs S}\le
\lambda^{\abs S}\le\lambda\).  So \(\abs H\le\lambda\); and the pairs
\((\{k\},\varnothing)\), \(k<\lambda\), show \(\abs H\ge\lambda\).  Fix an
enumeration \(H=\{(A_k,h_k):k<\lambda\}\).

\emph{The family.}  For \(B\subseteq\lambda\) put
\(f_B(k)=h_k(B\cap A_k)\) if \(B\cap A_k\) is in the domain of \(h_k\), and
\(f_B(k)=0\) otherwise.  Let \(F=\{f_B:B\subseteq\lambda\}\).

\emph{Distinct sets give distinct functions.}  If \(x\in B\setminus B'\), let
\(A=\{x\}\) and \(h(\{x\})=1\), \(h(\varnothing)=0\).  Then
\((A,h)=(A_k,h_k)\) for some \(k\), and \(f_B(k)=1\neq0=f_{B'}(k)\).  Hence
\(B\mapsto f_B\) is injective and \(\abs F=2^\lambda\).

\emph{\(\mu\)-consistency.}  Condition A) holds since
\(\{\lambda\}=[\{\lambda\}]\).  Since \(G=\varnothing\), condition B) concerns
only distinct \(f^i=f_{B_i}\in F\) and values \(j_i<\mu\), \(i<\chi<\mu\);
the \(B_i\) are distinct.  For each pair \(i_1\neq i_2\) choose a point of
\(B_{i_1}\mathbin\triangle B_{i_2}\), and let \(A\) be the set of chosen
points; \(\abs A\le\chi\cdot\chi<\mu\).  Then \(i_1\neq i_2\) implies
\(A\cap B_{i_1}\neq A\cap B_{i_2}\), so
\(S=\{A\cap B_i:i<\chi\}\) has \(\abs S=\chi<\mu\), and
\(h(A\cap B_i)=j_i\) defines a function \(h:S\to\mu\).  As \((A,h)\in H\),
\((A,h)=(A_k,h_k)\) for some \(k\), and
\(f^i(k)=h_k(B_i\cap A_k)=j_i\) for all \(i<\chi\).  So
\(\{k:f^i(k)=j_i\text{ for }i<\chi\}\) is nonempty, which is what
\(\neq\varnothing\ (\mathrm{mod}\ \{\lambda\})\) means, because
\(Z=\varnothing\ (\mathrm{mod}\ \{\lambda\})\) says
\(Z\subseteq\lambda-\lambda\).
\end{proof}
\end{annexitem}

\begin{annexitem}{annex:p03-w03-p03-017-lemma2}{thm:p03-w03-p03-017-lemma2}{Lemma~2: monotonicity and limits}
\begin{proof}
The source says only that the proof is immediate
\cite[p.~227]{Shelah1971Ultrapowers}; here are the steps.  Condition B) of
\cref{def:p03-w03-p03-017-consistent} does not mention \(\kappa\), and
condition A) becomes weaker as \(\kappa\) grows; this gives A).  For C),
condition A) is unchanged, and every instance of B) for \((F',G')\) is an
instance for \((F,G)\).

For B), fix a cofinal sequence \(\langle\delta_\xi:\xi<\mathrm{cf}(\delta)\rangle\)
in \(\delta\) and generating families \(E_i\) of \(D_i\) with
\(\abs{E_i}\le\kappa_i\).  The union \(D\) of the increasing chain of filters
\(D_i\) is a filter.  Let \(E=\bigcup_\xi E_{\delta_\xi}\); then
\(\abs E\le\mathrm{cf}(\delta)\cdot\kappa=\kappa\).  As \(E\subseteq D\),
\([E]\subseteq D\); and each \(D_i\) is contained in some
\(D_{\delta_\xi}=[E_{\delta_\xi}]\subseteq[E]\), so \(D=[E]\).  For condition
B), take an instance: \(f^i\in F\), \(j_i<\mu\) (\(i<\chi<\mu\)), and pairs
\(f^m\in F\), \(g^m\in G\) (\(m\le n\)), and suppose the displayed set \(Z\)
satisfies \(Z\subseteq\lambda-X\) with \(X\in D\).  Finitely many indices
\(i_0,\ldots,i_n,i'<\delta\) have \(g^m\in G_{i_m}\) and \(X\in D_{i'}\); let
\(i^*\) be their maximum.  Then \(g^m\in G_{i^*}\), \(X\in D_{i^*}\), and all
listed \(f\)'s lie in \(F\subseteq F_{i^*}\), so this is a violation of B) for
\((F_{i^*},G_{i^*},D_{i^*})\), which is impossible.
\end{proof}
\end{annexitem}

\begin{annexitem}{annex:p03-w03-p03-017-lemma3}{thm:p03-w03-p03-017-lemma3}{Lemma~3: bounded-range functions}
\begin{proof}
Following \cite[pp.~227--228]{Shelah1971Ultrapowers}, close a generating
family of \(D\), together with \(\lambda\), under finite intersections; the
result \(E\) has \(\abs E\le\kappa\) and \([E]=D\), so every \(X\in D\)
contains some member of \(E\).  Enumerate \(E=\{J_\alpha:\alpha<\kappa\}\)
(repetitions allowed).  Then \(Z=\varnothing\ (\mathrm{mod}\ D)\) iff
\(Z\subseteq\lambda-J_\alpha\) for some \(\alpha<\kappa\).

\emph{Reduction to finite \(G_1\).}  It suffices to find, for each finite
\(G_1\subseteq G\), a set \(F(G_1)\subseteq F\) with
\(\abs{F(G_1)}\le\kappa\) such that \((F-F(G_1),G_1,D)\) is
\(\kappa\)-consistent.  Indeed, put
\(F'=F-\bigcup\{F(G_1):G_1\subseteq G\text{ finite}\}\); there are at most
\(\kappa\) finite subsets of \(G\) since \(\abs G\le\kappa\), so
\(\abs{F-F'}\le\kappa\).  An instance of B) for \((F',G,D)\) uses finitely
many \(g\)'s, forming some finite \(G_1\), and functions in
\(F'\subseteq F-F(G_1)\); it is an instance for \((F-F(G_1),G_1,D)\) and so
cannot fail.

\emph{The finite case.}  Suppose that for some finite \(G_1\) no such
\(F(G_1)\) exists.  Call a \emph{violation} for a subset \(F^*\subseteq F\) a
system consisting of distinct functions \(f_i,f^*_m\in F^*\)
(\(i<\chi<\mu\), \(m\le n\)), values \(j_i<\mu\), a sequence
\(s=\langle g_0,\ldots,g_n\rangle\) of members of \(G_1\), and
\(\alpha<\kappa\), such that
\[
 \{k<\lambda:f_i(k)=j_i\ (i<\chi),\ f^*_m(k)=g_m(k)\ (m\le n)\}
 \subseteq\lambda-J_\alpha .
\]
By recursion on \(\beta<\kappa^+\) choose violations \(V_\beta\) for
\(F-U_\beta\), where \(U_\beta\) is the set of functions used by the
\(V_{\beta'}\), \(\beta'<\beta\).  This is possible: each violation uses
fewer than \(\mu\) functions, so \(\abs{U_\beta}\le\abs\beta\cdot\mu\le\kappa\);
by assumption \((F-U_\beta,G_1,D)\) is not \(\kappa\)-consistent; its
condition A) holds, so B) fails, and a failure of B) is a violation.  The
\(V_\beta\) use pairwise disjoint sets of functions.

The data \((\chi,\alpha,s)\) of a violation range over a set of size at most
\(\mu\cdot\kappa\cdot\aleph_0=\kappa\), so some value
\((\chi^0,\alpha^0,s^0)\), with \(s^0=\langle g_0,\ldots,g_n\rangle\), is taken
by \(\kappa^+\) of the \(V_\beta\) (otherwise \(\kappa^+\) would be a union of
at most \(\kappa\) sets each of size at most \(\kappa\)).  Each \(g_m\) maps
\(\lambda\) into a cardinal \(\chi(g_m)<\mu\); let \(\chi^*=\max_m\chi(g_m)<\mu\)
and let \(\Sigma\) be the set of sequences
\(\langle j^*_0,\ldots,j^*_n\rangle\) of ordinals below \(\chi^*\).  Then
\(\abs\Sigma=(\chi^*)^{n+1}<\mu\).  Assign to each \(\sigma\in\Sigma\) a
different violation \(V_{\beta(\sigma)}\) with data \((\chi^0,\alpha^0,s^0)\),
using functions \(f^\sigma_i\) (\(i<\chi^0\)) with values \(j^\sigma_i\) and
\(f^{*\sigma}_m\) (\(m\le n\)).  Let
\[
 A=\{k<\lambda:\text{for all }\sigma\in\Sigma,\ i<\chi^0,\ m\le n:\
 f^\sigma_i(k)=j^\sigma_i,\ f^{*\sigma}_m(k)=\sigma(m)\}.
\]
These are \(\abs\Sigma\cdot(\chi^0+n+1)<\mu\) distinct functions from \(F\)
with prescribed values below \(\mu\), so by condition B) for
\((F,\varnothing,D)\), \(A\neq\varnothing\ (\mathrm{mod}\ D)\); in particular
\(A\not\subseteq\lambda-J_{\alpha^0}\).  Pick \(k\in A\cap J_{\alpha^0}\) and
let \(\sigma=\langle g_0(k),\ldots,g_n(k)\rangle\in\Sigma\).  Then
\(f^\sigma_i(k)=j^\sigma_i\) and \(f^{*\sigma}_m(k)=g_m(k)\), so \(k\) lies in
the set of the violation \(V_{\beta(\sigma)}\), which is contained in
\(\lambda-J_{\alpha^0}\).  This contradicts \(k\in J_{\alpha^0}\).

The source indexes the sequences by ordinals below \(\chi^*\); only
\(\abs\Sigma<\mu\) is used here, which also covers finite \(\chi^*\).
\end{proof}
\end{annexitem}

\begin{annexitem}{annex:p03-w03-p03-017-lemma4}{thm:p03-w03-p03-017-lemma4}{Lemma~4: deciding sets}
\begin{proof}
\emph{A)} \cite[pp.~228--229]{Shelah1971Ultrapowers}.  Let
\(D_1=[D\cup\{A\}]\) and \(D_2=[D\cup\{\lambda-A\}]\).  If \(D=[E]\) with
\(\abs E\le\kappa\), then \(D_1=[E\cup\{A\}]\) and \(D_2=[E\cup\{\lambda-A\}]\),
so condition A) holds for both.  If \((F,G,D_1)\) is \(\kappa\)-consistent,
take \(F'=F\).  Otherwise B) fails for it: there are distinct
\(f_i,f^m\in F\), values \(j_i\), and \(g^m\in G\) (\(i<\chi<\mu\),
\(m\le n\)) whose set \(B\) lies in \(\lambda-Y\) for some \(Y\in D_1\).  Every
member of \(D_1\) contains \(X\cap A\) for some \(X\in D\), so
\(B\subseteq(\lambda-A)\cup(\lambda-X)\).  Let
\(F'=F-\{f_i,f^m:i<\chi,\ m\le n\}\); then \(\abs{F-F'}<\mu\).  If
\((F',G,D_2)\) is \(\kappa\)-consistent we are done.  Otherwise there are
distinct \(f^*_i,f^{*m}\in F'\), values \(j^*_i\) and \(g^{*m}\in G\)
(\(i<\chi^*<\mu\), \(m\le n^*\)) whose set \(B^*\) satisfies
\(B^*\subseteq(\lambda-X^*)\cup A\) for some \(X^*\in D\).  If \(k\in B\cap
B^*\) and \(k\in A\), then \(k\notin X\); if \(k\notin A\), then
\(k\notin X^*\).  So \(B\cap B^*\subseteq\lambda-(X\cap X^*)\) and
\(B\cap B^*=\varnothing\ (\mathrm{mod}\ D)\).  But \(B\cap B^*\) is the set
of the combined instance, whose \(f\)'s are distinct (the starred ones lie in
\(F'\)) and number \(\chi+\chi^*<\mu\) plus finitely many pairs; this
contradicts the \(\kappa\)-consistency of \((F,G,D)\).  When \(D_1\) or
\(D_2\) is improper the same argument applies, as the source notes.

\emph{B)}  By recursion on \(\alpha\le\kappa\) define
\(F^\alpha\), \(D^\alpha\) with \(F^0=F\), \(D^0=D\), \((F^\alpha,G,D^\alpha)\)
\(\kappa\)-consistent.  At \(\alpha+1\) apply A) to \(A_\alpha\), obtaining
\(F^{\alpha+1}\subseteq F^\alpha\) with
\(\abs{F^\alpha-F^{\alpha+1}}<\mu\) and \(D^{\alpha+1}\) equal to
\([D^\alpha\cup\{A_\alpha\}]\) or \([D^\alpha\cup\{\lambda-A_\alpha\}]\).  At
limits \(\delta\le\kappa\) take intersections and unions; Lemma~2B applies with
all \(\kappa_i=\kappa\), since \(\mathrm{cf}(\delta)\le\abs\delta\le\kappa\).
Put \(F'=F^\kappa\) and \(D'=D^\kappa\).  Then
\(\abs{F-F'}\le\kappa\cdot\mu=\kappa\) because \(\mu\le\kappa\), \(D'\supseteq D\)
is a filter, and \(A_\alpha\) or \(\lambda-A_\alpha\) lies in
\(D^{\alpha+1}\subseteq D'\).
\end{proof}
\end{annexitem}

\begin{annexitem}{annex:p03-w03-p03-017-lemma5}{thm:p03-w03-p03-017-lemma5}{Lemma~5: realizing a type by a new function}
\begin{proof}
Following \cite[pp.~229--230]{Shelah1971Ultrapowers}, enumerate
\(\dom M=\{c_j:j<\chi\}\); \(\chi\ge1\).  For \(l<l_0\) define
\(g_l:\lambda\to\chi\) by letting \(g_l(k)\) be the least \(j\) with
\(M\models\phi_l[c_j,\bar a_l[k]]\) if
\(M\models(\exists x)\phi_l(x,\bar a_l[k])\), and \(0\) otherwise.  Let
\(G=\{g_l:l<l_0\}\); \(\abs G\le\kappa\) and \(\mu+\abs G\le\kappa\).  By
Lemma~3 there is \(F_1\subseteq F\) with \(\abs{F-F_1}\le\kappa\) and
\((F_1,G,D)\) \(\kappa\)-consistent.  Since \(\abs F>\kappa\), \(F_1\neq\varnothing\);
choose \(f\in F_1\), let \(F'=F_1-\{f\}\), and put
\(\bar a[k]=c_{f(k)}\) if \(f(k)<\chi\) and \(\bar a[k]=c_0\) otherwise.  Let
\(A_l=\{k:M\models\phi_l[\bar a[k],\bar a_l[k]]\}\),
\(E=\{A_l:l<l_0\}\) and \(D'=[D\cup E]\).  Then \(\abs{F-F'}\le\kappa\),
\(D'\supseteq D\), and \(D'\) is generated by a generating family of \(D\)
together with \(E\), at most \(\kappa\) sets.  By the closure of \(p\) under
conjunctions, \(A_l\cap A_{l'}=A_{l''}\), so every member of \(D'\) contains
\(X\cap A_l\) for some \(X\in D\) and \(l<l_0\) (if \(l_0=0\) then \(D'=D\)
and there is nothing to prove).

Suppose condition B) fails for \((F',\varnothing,D')\): there are distinct
\(f_i\in F'\) and \(j_i<\mu\) (\(i<\chi_1<\mu\)) with
\(\bar A=\{k:f_i(k)=j_i\ (i<\chi_1)\}\subseteq\lambda-X'\), \(X'\in D'\).  Take
\(X\in D\), \(l<l_0\) with \(X'\supseteq X\cap A_l\); then
\(\bar A\cap A_l\cap X=\varnothing\).  Let
\(Y=\{k:f_i(k)=j_i\ (i<\chi_1),\ f(k)=g_l(k)\}\).  If \(k\in Y\cap X\cap A^l\),
then \(M\models(\exists x)\phi_l(x,\bar a_l[k])\), so
\(M\models\phi_l[c_{g_l(k)},\bar a_l[k]]\); as \(f(k)=g_l(k)<\chi\),
\(\bar a[k]=c_{g_l(k)}\), so \(k\in A_l\); and \(k\in\bar A\cap X\), which is
impossible.  Hence \(Y\subseteq\lambda-(X\cap A^l)\) with \(X\cap A^l\in D\),
that is, \(Y=\varnothing\ (\mathrm{mod}\ D)\).  The functions \(f_i\) and
\(f\) are distinct members of \(F_1\) and \(g_l\in G\), so this contradicts
the \(\kappa\)-consistency of \((F_1,G,D)\).  The case \(\chi_1=0\) shows that
\(D'\) is proper.  Finally \(A_l\in D'\) for every \(l\).
\end{proof}
\end{annexitem}

\begin{annexitem}{annex:p03-w03-p03-017-language}{thm:p03-w03-p03-017-language}{Reduction to a fixed language of cardinality \(\lambda\)}
\begin{proof}
For relation symbols of the same positive arity put \(R\sim S\) when
\(M\models\forall\bar x(R\bar x\leftrightarrow S\bar x)\); for function symbols
of the same arity \(f\sim g\) when \(M\models\forall\bar x\,f\bar x=g\bar x\); for
constants \(c\sim d\) when \(M\models c=d\).  As \(M\equiv N\), the same
sentences hold in \(N\).  Choose one representative in each class, and let
\(L'\) be the set of representatives.  Write \(\nu=\lVert M\rVert<\mu\).  The
classes of \(n\)-ary relation symbols inject into \(\mathcal P((\dom M)^n)\),
of size \(2^{\nu^n}\le\lambda^{\nu^n}\le\lambda\) since \(\nu^n<\mu\); those of
function symbols into \((\dom M)^{(\dom M)^n}\), of size at most
\(\lambda^{\nu^n}\le\lambda\); constants into \(\dom M\).  So there is an
arity-preserving injection of \(L'\) into \(L^0\); rename \(L'\) along it to
\(L_*\), and transport \(M\restriction L'\) and \(N\restriction L'\) along
bijections of their universes with the cardinals \(\lVert M\rVert\),
\(\lVert N\rVert\).  This gives \(M_*\) and \(N_*\); they are elementarily
equivalent because \(L'\)-sentences are \(L\)-sentences and transport and
renaming preserve truth (\cref{lem:iso-invariance}).  Nullary relation
symbols of \(L\) are not represented; each has the same truth value in
\(M\) and \(N\) because \(M\equiv N\).

Let \(D\) be an ultrafilter over \(\lambda\).  Ultrapowers are computed symbol
by symbol, so \((M\restriction L')^\lambda/D\) is the \(L'\)-reduct of
\(M^\lambda/D\), and transport and renaming carry
\((M\restriction L')^\lambda/D\) onto \(M_*^\lambda/D\) (item~1 of
\cref{thm:p03-w03-p03-018-comparison}); the same computation with \(N\),
\(N\restriction L'\) and \(N_*\) in place of \(M\), \(M\restriction L'\) and
\(M_*\) carries \((N\restriction L')^\lambda/D\) onto \(N_*^\lambda/D\).  By
\cref{cor:diagonal-ultrapower}, \(M^\lambda/D\) satisfies every sentence
true in \(M\), in particular \(\forall\bar x(S\bar x\leftrightarrow S'\bar x)\)
and its analogues for each symbol \(S\) and its representative \(S'\), and
the truth value of each nullary symbol; likewise for \(N^\lambda/D\).

\emph{Isomorphism.}  An isomorphism \(M_*^\lambda/D\cong N_*^\lambda/D\) gives
an \(L'\)-isomorphism \(\Phi\) of the \(L'\)-reducts of \(M^\lambda/D\) and
\(N^\lambda/D\).  For any \(S\in L\) with representative \(S'\), \(S\) and
\(S'\) have the same interpretation in \(M^\lambda/D\) and in \(N^\lambda/D\),
so \(\Phi\) preserves \(S\); a nullary symbol has the same value on both
sides.  So \(\Phi\) is an \(L\)-isomorphism.

\emph{Types.}  Replace in an \(L\)-formula every symbol by its representative
and every nullary relation symbol by \(\forall x\,x=x\) or its negation
according to its value in \(M\).  The result is an \(L'\)-formula equivalent
to the original in \(M^\lambda/D\) for all parameters, by the sentences just
listed.  A set of \(L\)-formulas with at most \(\kappa\) parameters, finitely
satisfiable in \(M^\lambda/D\), translates into a set of \(L'\)-formulas with
the same parameters and the same satisfying tuples; transported to
\(M_*^\lambda/D\) it is finitely satisfiable there, hence realized, and the
transport back of a realizing tuple realizes the original set.

\emph{Counting.}  The sublanguages \(L_*\subseteq L^0\) number \(2^\lambda\).
For a cardinal \(\chi<\mu\), each symbol of \(L_*\) has at most
\(\max(2^{\chi^n},\chi^{\chi^n},\chi)\le\lambda\) interpretations on \(\chi\),
so there are at most \(\lambda^\lambda=2^\lambda\) \(L_*\)-structures with
universe \(\chi\).  With at most \(\mu\le\lambda\) cardinals \(\chi\), the pairs
number at most \(2^\lambda\) \cite[p.~230]{Shelah1971Ultrapowers}.
\end{proof}
\end{annexitem}

\begin{annexitem}{annex:p03-w03-p03-017-bookkeeping}{thm:p03-w03-p03-017-bookkeeping}{The bookkeeping function \(R\)}
\begin{proof}
The source asserts the existence of \(R\) from cardinality considerations
\cite[p.~230]{Shelah1971Ultrapowers}; here they are.  Regard the four kinds
of value as tagged, so that they are pairwise distinct.  Let \(V_{\mathrm A}\),
\(V_{\mathrm B}\), \(V_{\mathrm C}\), \(V_{\mathrm D}\) be the sets of all values
required in A), B), C), D).

\emph{Sizes.}  \(\abs{V_{\mathrm A}},\abs{V_{\mathrm B}}\le\theta\cdot\theta=\theta\)
and \(\abs{V_{\mathrm D}}=\theta\).  Each \(L_i\) has at most \(\lambda\)
symbols, so there are at most \(\theta\) formulas in the variables \(x\),
\(y_\alpha\) (\(\alpha<\theta\)).  For a cardinal \(\nu\ge1\) with
\(2^\nu\le\theta\), the sets \(p\) of size \(\nu\) number at most
\(\theta^\nu=2^{\lambda\cdot\nu}=\max(2^\lambda,2^\nu)=\theta\).  There are at
most \(\theta\) such \(\nu\), so \(\abs{V_{\mathrm C}}\le\theta\cdot\theta\cdot
\theta=\theta\).

\emph{A cofinality bound.}  If \(2^\nu\le\theta\) then
\(\nu<\mathrm{cf}(\theta)\).  Indeed, write \(\theta\) as the union of a
cofinal sequence \(\langle\theta_\xi:\xi<\mathrm{cf}(\theta)\rangle\) of
ordinals below \(\theta\); by K\"onig's theorem
(\cref{thm:p06-konig-cardinal}),
\(\theta\le\sum_\xi\abs{\theta_\xi}<\prod_\xi\theta=\theta^{\mathrm{cf}(\theta)}\).
If \(\nu\ge\mathrm{cf}(\theta)\) we would get
\(\theta<\theta^{\mathrm{cf}(\theta)}\le\theta^\nu=2^{\lambda\cdot\nu}=
\max(\theta,2^\nu)=\theta\).  Applied to \(\nu=\lambda\) and \(\nu=\abs p\),
this shows that each \(p\) in \(V_{\mathrm C}\) has a set \(P(p)\) of parameter
indices with \(\abs{P(p)}\le\abs p+\lambda<\mathrm{cf}(\theta)\).

\emph{Positions.}  Each \(\gamma<\theta\) is uniquely \(\delta+n\) with
\(\delta\) zero or a limit and \(n<\omega\); call \(\gamma\) even or odd
according to \(n\).  There are \(\theta\) even and \(\theta\) odd ordinals
below \(\theta\).  Fix a bijection \(e\) from the even ordinals onto
\(V_{\mathrm A}\cup V_{\mathrm B}\cup V_{\mathrm D}\) and let \(R(\gamma)=e(\gamma)\)
for even \(\gamma\).  For \(c=\langle i,p\rangle\in V_{\mathrm C}\) let \(b(c)\)
be the supremum of \(e^{-1}(\langle i,1,\bar a^i_\alpha\rangle)+1\) over
\(\alpha\in P(p)\); by the cofinality bound \(b(c)<\theta\).  Enumerate
\(V_{\mathrm C}=\{c_\zeta:\zeta<\theta\}\), repetitions allowed, and choose by
recursion odd ordinals \(o(\zeta)>b(c_\zeta)\) different from all
\(o(\zeta')\), \(\zeta'<\zeta\); this is possible because there are
\(\theta\) odd ordinals above \(b(c_\zeta)\) and fewer than \(\theta\) have
been used.  Put \(R(o(\zeta))=c_\zeta\), and \(R(\gamma)=\langle0,
\varnothing\rangle\) at the remaining odd \(\gamma\); the empty set of
formulas is closed under conjunctions and has no parameters.

A), B), D) hold because \(e\) is onto; C) holds with \(\gamma=o(\zeta)\) and
\(\gamma_\alpha=e^{-1}(\langle i,1,\bar a^i_\alpha\rangle)<b(c_\zeta)<\gamma\),
and the same inequality gives the clause about every \(\gamma\) with
\(R(\gamma)=\langle i,p\rangle\); E) holds since every value has one kind.
\end{proof}
\end{annexitem}

\begin{annexitem}{annex:p03-w03-p03-017-induction}{thm:p03-w03-p03-017-induction}{The inductive construction}
\begin{proof}
The recursion is on \(\gamma\le\theta\), with the choices made by a fixed
choice function.  Write \(\kappa_\gamma=\lambda+\abs\gamma\).  Two remarks
are used throughout.  First, for \(\gamma<\theta\), condition~1 gives
\(\abs{F_\gamma}=\theta>\kappa_\gamma\), since \(\abs{F_0}=\theta\) and
\(\abs{F_0-F_\gamma}\le\kappa_\gamma<\theta\).  Second, conditions~3 and~4
together decide formulas on the \(N_i\)-side too: if
\((\bar a_m,\bar b_m)\in H^i_\gamma\), the set
\(\{k:N_i\models\phi[\bar b_1[k],\ldots]\}\) or its complement contains the
intersection of a member of \(D_\gamma\) given by~4 with the member given
by~3.

\emph{Stage \(0\)} \cite[p.~232, Case~I]{Shelah1971Ultrapowers}.  Let
\(F_0\) be the family of Lemma~1, \(D_0=\{\lambda\}\), \(H^i_0=\varnothing\).
By Lemma~2A, \((F_0,\varnothing,D_0)\) is \(\lambda\)-consistent, since
\(\mu\le\lambda\).  Conditions~3 and~4 concern sentences only: as
\(M_i\equiv N_i\), the set in~3 is \(\lambda\), and in~4 one of the two sets
is \(\lambda\).

\emph{Limit \(\gamma\)} \cite[p.~232, Case~II]{Shelah1971Ultrapowers}.  Put
\(F_\gamma=\bigcap_{\beta<\gamma}F_\beta\),
\(D_\gamma=\bigcup_{\beta<\gamma}D_\beta\),
\(H^i_\gamma=\bigcup_{\beta<\gamma}H^i_\beta\).  Lemma~2B with
\(\kappa=\kappa_\gamma\ge\kappa_\beta\) and \(\kappa_\gamma\ge\abs\gamma\ge
\mathrm{cf}(\gamma)\) gives the consistency in~1; further
\(\abs{F_0-F_\gamma}\le\abs\gamma\cdot\kappa_\gamma=\kappa_\gamma\) and
\(\abs{\bigcup_iH^i_\gamma}\le\abs\gamma\cdot\abs\gamma=\abs\gamma\).  Finitely
many pairs or elements of \(H^i_\gamma\) already lie in some \(H^i_\beta\),
so~3 and~4 hold with \(D_\beta\subseteq D_\gamma\).  Conditions~5--8 concern
successor stages.

\emph{Successor \(\gamma=\beta+1\).}  Let \(\kappa=\kappa_\beta\),
\(F=F_\beta\), \(D=D_\beta\), \(H^i=H^i_\beta\).  In each case below we
produce \(F'\subseteq F\) with \(\abs{F-F'}\le\kappa\), a filter
\(D'\supseteq D\) with \((F',\varnothing,D')\) \(\kappa\)-consistent, and
\(H^{i'}_\gamma\supseteq H^{i'}\) with at most one new pair in all, and set
\(F_\gamma=F'\), \(D_\gamma=D'\).  Then~1 holds at \(\gamma\) by Lemma~2A and
\(\abs{F_0-F_\gamma}\le\kappa+\kappa\le\kappa_\gamma\); \(2\) holds since
\(\abs\beta+1=\abs\gamma\); and~3, \(4\) hold for tuples from \(H^{i'}\)
because \(D\subseteq D'\).  When a case requires nothing, everything is left
unchanged.

\emph{Case III: \(R(\beta)=\langle i,1,\bar a\rangle\)} \cite[p.~232]{Shelah1971Ultrapowers}.
If \(\bar a\in\mathrm{Dom}\,H^i\), leave everything unchanged.  Otherwise let
\(\Phi\) be the set of pairs \((\psi,s)\) with \(\psi(x,x_1,\ldots,x_n)\) an
\(L_i\)-formula and \(s=\langle(\bar a_1,\bar b_1),\ldots,(\bar a_n,\bar
b_n)\rangle\) a finite sequence from \(H^i\); \(\abs\Phi\le\lambda+\abs\beta=
\kappa\).  Put
\(A_{\psi,s}=\{k:M_i\models\psi[\bar a[k],\bar a_1[k],\ldots,\bar a_n[k]]\}\).
Lemma~4B (with \(\mu\le\kappa\)) gives \(F^1\), \(D^1\supseteq D\), with
\(\abs{F-F^1}\le\kappa\), \((F^1,\varnothing,D^1)\) \(\kappa\)-consistent, and
every \(A_{\psi,s}\) or its complement in \(D^1\).  Let
\(p=\{(\psi,s)\in\Phi:A_{\psi,s}\in D^1\}\), read in \(N_i\) as the formulas
\(\psi(x,\bar b_1,\ldots,\bar b_n)\).  It is closed under conjunctions: for
\((\psi,s),(\psi',s')\in p\) the conjunction, with the variables of \(\psi'\)
shifted past those of \(\psi\) and the sequence \(s^\frown s'\), has set
\(A_{\psi,s}\cap A_{\psi',s'}\in D^1\).  For \((\psi,s)\in p\),
\(A_{\psi,s}\subseteq\{k:M_i\models\exists x\,\psi(x,\bar a_s[k])\}\), and
condition~3 at \(\beta\) for \(\exists x\,\psi\) gives a member of \(D\) on
which this \(M_i\)-set agrees with
\(\{k:N_i\models\exists x\,\psi(x,\bar b_s[k])\}\); so the latter is in
\(D^1\).  Lemma~5, applied to \(N_i\) (of cardinality \({<}\,\mu\)), \(p\),
\((F^1,\varnothing,D^1)\), with \(\abs{F^1}=\theta>\kappa\), gives
\(\bar b\in(\dom N_i)^\lambda\), \(F^2\subseteq F^1\) with
\(\abs{F^1-F^2}\le\kappa\), and \(D^2\supseteq D^1\) with
\((F^2,\varnothing,D^2)\) \(\kappa\)-consistent and
\(B_{\psi,s}=\{k:N_i\models\psi[\bar b[k],\bar b_s[k]]\}\in D^2\) for
\((\psi,s)\in p\).  Put \(F'=F^2\), \(D'=D^2\), \(H^i_\gamma=H^i\cup\{(\bar
a,\bar b)\}\).  Condition~5 holds.  For~3 at \(\gamma\), a formula applied to
a list of pairs containing \((\bar a,\bar b)\) is \(\psi(x,\ldots)\) applied
to that pair at \(x\) and to a sequence \(s\) of earlier pairs.  If
\(A_{\psi,s}\in D^1\) then \((\psi,s)\in p\), and \(A_{\psi,s}\cap B_{\psi,s}
\in D^2\) lies in the set of~3.  If \(\lambda-A_{\psi,s}\in D^1\) then
\((\neg\psi,s)\in p\), so \(\lambda-B_{\psi,s}\in D^2\), and
\((\lambda-A_{\psi,s})\cap(\lambda-B_{\psi,s})\) lies in the set of~3.
Condition~4 at \(\gamma\) holds by the decisions in \(D^1\).

\emph{Case IV: \(R(\beta)=\langle i,2,\bar b\rangle\)}.  The source says this
is the same as Case~III \cite[p.~232]{Shelah1971Ultrapowers}.  If
\(\bar b\in\mathrm{Range}\,H^i\), leave everything unchanged.  Otherwise
decide the sets \(B_{\psi,s}=\{k:N_i\models\psi[\bar b[k],\bar b_s[k]]\}\) by
Lemma~4B, form \(p\) from those in \(D^1\), read in \(M_i\) with parameters
\(\bar a_s\); its existential sets are in \(D^1\) by condition~3; apply
Lemma~5 to \(M_i\) to obtain \(\bar a\) with
\(A_{\psi,s}\in D^2\) for \((\psi,s)\in p\), and add the pair
\((\bar a,\bar b)\).  The argument for~3 is the previous one with \(M_i\)
and \(N_i\) exchanged.  For~4, a tuple from \(\mathrm{Dom}\,H^i_\gamma\)
containing \(\bar a\) gives some \((\psi,s)\); either \((\psi,s)\in p\), so
\(A_{\psi,s}\in D^2\), or \((\neg\psi,s)\in p\), so \(\lambda-A_{\psi,s}\in
D^2\).  This holds even if \(\bar a\) already occurs in an earlier pair,
which is why \(H^i_\gamma\) is a set of pairs.  Condition~6 holds.

\emph{Case V: \(R(\beta)=\langle i,p\rangle\)} \cite[p.~232]{Shelah1971Ultrapowers}.
If the hypothesis of~7 fails, leave everything unchanged.  Otherwise, by
\cref{thm:p03-w03-p03-017-bookkeeping} each parameter function of \(p\) is
\(R(\gamma')\) for some \(\gamma'<\beta\), so the distinct parameter
functions number at most \(\abs\beta\).  Identify two variables \(y_\alpha\),
\(y_{\alpha'}\) when \(\bar a^i_\alpha=\bar a^i_{\alpha'}\); this turns \(p\)
into a set \(\hat p\), closed under conjunctions, of at most
\(\lambda+\abs\beta=\kappa\) formulas, with the same sets as \(p\).  Lemma~5
applied to \(M_i\), \(\hat p\) and \((F,\varnothing,D)\) gives \(\bar a\),
\(F'\), \(D'\) with condition~7.

\emph{Case VI: \(R(\beta)=A\subseteq\lambda\)} \cite[p.~232]{Shelah1971Ultrapowers}.
Lemma~4A gives \(F'\) with \(\abs{F-F'}<\mu\le\kappa\) and \(D'\) equal to
\([D\cup\{A\}]\) or \([D\cup\{\lambda-A\}]\), with \((F',\varnothing,D')\)
\(\kappa\)-consistent.  Condition~8 holds.

By E) of \cref{thm:p03-w03-p03-017-bookkeeping} exactly one case applies at
each successor, so the recursion is defined up to \(\theta\).
\end{proof}
\end{annexitem}

\begin{annexitem}{annex:p03-w03-p03-017-shelah-proof}{thm:p03-w03-p03-017-shelah-proof}{Parts~1 and~2 of Shelah's theorem}
\begin{proof}
\emph{\(D\) is an ultrafilter} \cite[p.~231]{Shelah1971Ultrapowers}.  By
condition~1, \((F_\theta,\varnothing,D)\) is consistent, and the empty
instance of B) shows \(\varnothing\notin D\).  For \(A\subseteq\lambda\),
D) gives \(\gamma\) with \(R(\gamma)=A\), and~8 puts \(A\) or \(\lambda-A\)
into \(D_{\gamma+1}\subseteq D\).

\emph{Part 1 for listed pairs.}  Let \(H^i=H^i_\theta\).  By~5 and A),
\(\mathrm{Dom}\,H^i=(\dom M_i)^\lambda\); by~6 and B),
\(\mathrm{Range}\,H^i=(\dom N_i)^\lambda\).  Let
\(\hat H=\{(\bar a/D,\bar b/D):(\bar a,\bar b)\in H^i\}\).  For pairs
\((\bar a,\bar b),(\bar a',\bar b')\in H^i\), condition~3 for \(x_1=x_2\)
gives \(Z\in D\) on which \(\bar a[k]=\bar a'[k]\) iff \(\bar b[k]=\bar
b'[k]\); intersecting with \(Z\) shows that \(\bar a/D=\bar a'/D\) iff \(\bar
b/D=\bar b'/D\).  So \(\hat H\) is a bijection of \(M_i^\lambda/D\) onto
\(N_i^\lambda/D\).  For a relation symbol \(R\) and pairs
\((\bar a_m,\bar b_m)\), condition~3 gives a member of \(D\) on which the
\(M_i\)-set and the \(N_i\)-set of \(R\) agree; so one is in \(D\) iff the
other is, and by the definition of the ultrapower \(\hat H\) preserves
\(R\).  For a function symbol \(f\), let \(\bar c\) be the coordinatewise
value \(f(\bar a_1,\ldots,\bar a_n)\) and \((\bar c,\bar d)\in H^i\); the
formula \(f(x_1,\ldots,x_n)=x_{n+1}\) has \(M_i\)-set \(\lambda\), so its
\(N_i\)-set is in \(D\), that is,
\(f(\bar b_1/D,\ldots,\bar b_n/D)=\bar d/D=\hat H(\bar c/D)\) by
\cref{lem:los-terms}.  Constants are the case \(n=0\).  So \(\hat H\) is an
isomorphism \cite[p.~231]{Shelah1971Ultrapowers}.

\emph{Part 1 in general.}  Let \(M\equiv N\) have power \({<}\,\mu\).  By
\cref{thm:p03-w03-p03-017-language} there is a pair \((M_*,N_*)\) of the kind
listed, so \((M_*,N_*)=(M_i,N_i)\) for some \(i\), and
\(M_*^\lambda/D\cong N_*^\lambda/D\) yields \(M^\lambda/D\cong N^\lambda/D\).

\emph{Part 2.}  Let \(M\) have power \({<}\,\mu\), and let \(\kappa\) be
infinite with \(2^\kappa\le2^\lambda\).  Apply
\cref{thm:p03-w03-p03-017-language} to the pair \((M,M)\), and let
\((M_*,M_*)=(M_i,N_i)\).  We show that every set of formulas in finitely
many variables, with at most \(\kappa\) parameters, finitely satisfiable in
\(M_i^\lambda/D\), is realized there.

\emph{One variable.}  Let \(q(x)\) be such a set.  Each parameter is
\(\bar a^i_\alpha/D\) for some \(\alpha\); choose one such \(\alpha\) for each
parameter, write \(q\) with variables \(y_\alpha\), and let \(p\) be its
closure under finite conjunctions.  Then \(\abs p\le\lambda+\kappa\) and
\(2^{\abs p}\le2^\lambda\cdot2^\kappa=2^\lambda\).  Every \(\phi\in p\) is
satisfiable in \(M_i^\lambda/D\), so by \cref{thm:los}
\(S_\phi=\{k:M_i\models(\exists x)\phi(x,\bar a[k])\}\in D\), where
\(\bar a\) lists its parameter functions.  By C) there are \(\gamma\) with
\(R(\gamma)=\langle i,p\rangle\) and \(\gamma_\alpha<\gamma\) with
\(R(\gamma_\alpha)=\langle i,1,\bar a^i_\alpha\rangle\); by~5,
\(\bar a^i_\alpha\in\mathrm{Dom}\,H^i_{\gamma_\alpha+1}\subseteq
\mathrm{Dom}\,H^i_\gamma\).  By~4 at \(\gamma\), \(S_\phi\in D_\gamma\) or
\(\lambda-S_\phi\in D_\gamma\); the second would give
\(\varnothing\in D\).  So the hypothesis of~7 holds, and~7 gives
\(\bar a\in(\dom M_i)^\lambda\) whose sets lie in \(D_{\gamma+1}\subseteq D\)
for all \(\phi\in p\).  By \cref{thm:los}, \(\bar a/D\) realizes \(q\)
\cite[p.~231]{Shelah1971Ultrapowers}.

\emph{Finitely many variables.}  Induct on the number of variables.  Let
\(\Sigma(\bar x,x')\) be finitely satisfiable with at most \(\kappa\)
parameters; close it under finite conjunctions.  The set
\(\{\exists x'\sigma:\sigma\in\Sigma\}\) is finitely satisfiable (a finite
subset is implied by the existential of one conjunction in \(\Sigma\)), so it
is realized by some \(\bar c\).  Then \(\Sigma(\bar c,x')\) has at most
\(\kappa+\abs{\bar c}=\kappa\) parameters and is finitely satisfiable, so it
is realized by some \(d\), and \((\bar c,d)\) realizes \(\Sigma\).

\emph{Saturation.}  By \cref{thm:p03-w03-p03-017-language} the property just
proved passes to \(M^\lambda/D\).  A complete type over \(A\) with
\(\abs A<\kappa^+\) that is consistent with \(M^\lambda/D\) is finitely
satisfiable in it: if a conjunction \(\sigma\) of finitely many of its
formulas had no solution, \(\neg\exists\bar x\,\sigma\) would be true in the
structure with names for \(A\), contradicting consistency.  So it is
realized, and \(M^\lambda/D\) is \(\kappa^+\)-saturated in the sense of the
definition preceding \cref{thm:saturated-extension}.
\end{proof}
\end{annexitem}

\begin{annexitem}{annex:p03-w03-p03-017-keisler-shelah}{thm:p03-w03-p03-017-keisler-shelah}{Keisler--Shelah, both directions}
\begin{proof}
\emph{Forward.}  Let \(M\equiv N\), let
\(\nu=\lVert M\rVert+\lVert N\rVert+\aleph_0\) and \(\lambda=2^\nu\).  Then
\(\lambda^\nu=2^{\nu\cdot\nu}=\lambda\).  If \(\mu\le\nu\), monotonicity in
the exponent would give \(\lambda<\lambda^\mu\le\lambda^\nu=\lambda\); so
\(\nu<\mu\), and both models have power \({<}\,\mu\).  More generally any
infinite \(\lambda\) with \(\lVert M\rVert,\lVert N\rVert<\mu\) serves.  By
part~1 of \cref{thm:p03-w03-p03-017-shelah-proof},
\(M^\lambda/D\cong N^\lambda/D\) for the ultrafilter \(D\) constructed there;
take \(I=J=\lambda\) and \(U=V=D\).  No property of \(D\) beyond that theorem
is used, and the existential statement does not claim that a given index set
or ultrafilter works.

\emph{Converse.}  Let \(\Psi:M^I/U\cong N^J/V\), with \(U\), \(V\)
ultrafilters on nonempty sets.  For a sentence \(\sigma\): \(M\models\sigma\)
iff \(M^I/U\models\sigma\), by \L o\'s's theorem (\cref{thm:los}), because
\(\{i:M\models\sigma\}\) is \(I\in U\) or \(\varnothing\notin U\) (this is
\cref{cor:diagonal-ultrapower} for sentences); \(M^I/U\models\sigma\) iff
\(N^J/V\models\sigma\) by \cref{lem:iso-invariance}; and
\(N^J/V\models\sigma\) iff \(N\models\sigma\) as before.  So \(M\equiv N\).
This is the argument of item~2 of \cref{thm:p03-w03-p03-018-comparison}.
\end{proof}
\end{annexitem}

\subsection{The bound, saturated models and Keisler's characterizations}

\begin{annexitem}{annex:p03-w03-p03-017-mu-regular}{thm:p03-w03-p03-017-mu-regular}{The bound \(\mu\) is regular}
\begin{proof}
By \cref{thm:p03-w03-p03-018-mu}, \(\mu\) is an infinite cardinal, and
\(\lambda^\nu=\lambda\) for every cardinal \(\nu\) with \(1\le\nu<\mu\).
Suppose, for a contradiction, that \(\operatorname{cf}(\mu)<\mu\), and let
\(C\subseteq\mu\) be cofinal in \(\mu\) with order type
\(\operatorname{cf}(\mu)\) (\cref{def:p06-cofinality-regular-singular}).
Put \(\nu=\abs C\); then \(\nu\le\operatorname{cf}(\mu)<\mu\), and
\(\nu\ge1\) because \(\mu\neq0\).

An infinite cardinal is a limit ordinal, so \(\gamma+1<\mu\) for every
\(\gamma\in C\).  Every \(\beta<\mu\) satisfies \(\beta\le\gamma\) for some
\(\gamma\in C\), that is \(\beta\in\gamma+1\); hence
\(\mu=\bigcup_{\gamma\in C}(\gamma+1)\).  Consequently the map sending
\(f\in\lambda^\mu\) (a function from \(\mu\) to \(\lambda\)) to the family
\(\langle f\restriction(\gamma+1):\gamma\in C\rangle\) is injective, from
\(\lambda^\mu\) into \(\prod_{\gamma\in C}\lambda^{\gamma+1}\).

For \(\gamma\in C\) the cardinal \(\abs{\gamma+1}\) is at least \(1\) and
below \(\mu\), because \(\gamma+1<\mu\) and \(\mu\) is a cardinal; so the
set \(\lambda^{\gamma+1}\) has cardinality
\(\lambda^{\abs{\gamma+1}}=\lambda\).  Choose, by the Axiom of Choice, an
injection \(e_\gamma:\lambda^{\gamma+1}\to\lambda\) for each \(\gamma\in C\).
Then \(\langle x_\gamma\rangle_{\gamma\in C}\mapsto
\langle e_\gamma(x_\gamma)\rangle_{\gamma\in C}\) injects
\(\prod_{\gamma\in C}\lambda^{\gamma+1}\) into the set of functions from
\(C\) to \(\lambda\), whose cardinality is \(\lambda^\nu=\lambda\) since
\(1\le\nu<\mu\).  Composing the two injections gives
\(\lambda^\mu\le\lambda\), contrary to \(\lambda^\mu>\lambda\).  Hence
\(\operatorname{cf}(\mu)=\mu\), and the infinite cardinal \(\mu\) is
regular.
\end{proof}
\end{annexitem}

\begin{annexitem}{annex:p03-w03-p03-017-saturated-unique}{thm:p03-w03-p03-017-saturated-unique}{Uniqueness of saturated models of one power}
\begin{proof}
Call a map \(f\) from a subset of \(\dom M\) into \(\dom N\)
\emph{elementary} if \(M\models\varphi[\bar a]\) iff
\(N\models\varphi[f\bar a]\) for every formula \(\varphi(\bar x)\) and every
tuple \(\bar a\) from the domain of \(f\).  The empty map is elementary
because \(M\equiv N\).  An elementary map is injective (take \(\varphi\) to
be \(x_1=x_2\)), and the union of a chain of elementary maps is elementary,
since a formula involves finitely many elements, all in the domain of one
member of the chain.

\emph{Forth.}  Let \(f\) be elementary with domain \(X\), \(\abs X<\mu\), and
let \(m\in\dom M\).  Let \(p\) be the set of formulas
\(\varphi(x,f\bar a)\), with parameters from \(f[X]\), such that
\(M\models\varphi[m,\bar a]\).  For each formula \(\psi(x,f\bar a)\) exactly
one of \(\psi\) and \(\neg\psi\) is in \(p\), so \(p\) is complete.  If
\(\varphi_1(x,f\bar a_1),\ldots,\varphi_n(x,f\bar a_n)\in p\), then
\(M\models\exists x\bigwedge_k\varphi_k[x,\bar a_1,\ldots,\bar a_n]\), so
\(N\models\exists x\bigwedge_k\varphi_k[x,f\bar a_1,\ldots,f\bar a_n]\)
because \(f\) is elementary; thus every finite subset of \(p\) is realized
in \(N\), and \(p\) is consistent with \(N\).  Since \(\abs{f[X]}<\mu\) and
\(N\) is \(\mu\)-saturated, some \(n\in\dom N\) realizes \(p\), and then
\(f\cup\{(m,n)\}\) is elementary: a formula \(\varphi(x,\bar y)\) with
\(M\models\varphi[m,\bar a]\) puts \(\varphi(x,f\bar a)\) in \(p\), and one
with \(M\models\neg\varphi[m,\bar a]\) puts \(\neg\varphi(x,f\bar a)\) in
\(p\).

\emph{Back.}  The inverse of an elementary map is elementary from a subset
of \(\dom N\) into \(\dom M\); applying the forth step to it, with the roles
of \(M\) and \(N\) exchanged and the \(\mu\)-saturation of \(M\), adds any
given \(n\in\dom N\) to the range.

\emph{Construction.}  Enumerate \(\dom M=\{m_\zeta:\zeta<\mu\}\) and
\(\dom N=\{n_\zeta:\zeta<\mu\}\) (Choice).  Put \(f_0=\varnothing\); given
\(f_\zeta\) with \(\abs{f_\zeta}<\mu\), apply the forth step to add
\(m_\zeta\) to the domain and then the back step to add \(n_\zeta\) to the
range, obtaining \(f_{\zeta+1}\); at limits take unions.  By induction
\(\abs{f_\zeta}\le2\cdot\abs\zeta\), which is finite for finite \(\zeta\)
and equal to \(\abs\zeta<\mu\) otherwise; so \(\abs{f_\zeta}<\mu\) at each
stage \(\zeta<\mu\), as the steps require.  The union \(f\) of all
\(f_\zeta\) is elementary, has domain \(\dom M\) and range \(\dom N\).  An
elementary bijection preserves every atomic formula in both directions,
including those built from function and constant symbols, so it is an
isomorphism.
\end{proof}
\end{annexitem}

\begin{annexitem}{annex:p03-w03-p03-017-decomposition}{thm:p03-w03-p03-017-decomposition}{How the stages give Keisler's theorem}
\begin{proof}
If \(M\) is finite, then \(M\cong N\) by item~4 of
\cref{thm:p03-w03-p03-018-comparison}, and \(M^\lambda/D\cong N^\lambda/D\)
by item~1 there.  Let \(M\) be infinite; then so is \(N\), since for each
\(n\) the sentence saying that there are at least \(n\) elements is true in
\(M\).

\emph{Size.}  \(\abs{(\dom M)^\lambda}\le(\lambda^+)^\lambda
=(2^\lambda)^\lambda=2^{\lambda\cdot\lambda}=2^\lambda=\lambda^+\), and
\(M^\lambda/D\) is a quotient of \((\dom M)^\lambda\), so
\(\lVert M^\lambda/D\rVert\le\lambda^+\).  The diagonal map embeds \(M\)
into \(M^\lambda/D\) (\cref{cor:diagonal-ultrapower}), so \(M^\lambda/D\) is
infinite.  An infinite \(\lambda^+\)-saturated structure \(C\) has at least
\(\lambda^+\) elements: otherwise \(\abs{\dom C}\le\lambda<\lambda^+\), and
the set \(\{x\ne c:c\in\dom C\}\) is finitely realized in \(C\) because
\(C\) is infinite; by compactness
(\cref{thm:p03-compactness-ultraproduct}) it is realized by some element of
an elementary extension of \(C\), whose complete type over \(\dom C\) is
consistent with \(C\); by saturation that type is realized by some
\(c'\in\dom C\), which then satisfies \(x\ne c'\), a contradiction.  Hence
\(\lVert M^\lambda/D\rVert=\lambda^+\), and likewise
\(\lVert N^\lambda/D\rVert=\lambda^+\).

\emph{Isomorphism.}  \(M^\lambda/D\equiv M\equiv N\equiv N^\lambda/D\) by
\L o\'s's theorem (\cref{thm:los}, through \cref{cor:diagonal-ultrapower}).
By hypothesis both ultrapowers are \(\lambda^+\)-saturated, and both have
cardinality \(\lambda^+\).
\Cref{thm:p03-w03-p03-017-saturated-unique} with \(\lambda^+\) in place of
its cardinal gives \(M^\lambda/D\cong N^\lambda/D\).
\end{proof}
\end{annexitem}

\begin{annexitem}{annex:p03-w03-p03-017-gch-free-closed}{thm:p03-w03-p03-017-gch-free-closed}{Keisler's Theorems~2.6, 2.11 and~2.12 and Corollary~2.7 without the continuum hypothesis}
\begin{proof}
Write \(\mathfrak A\equiv\mathfrak B\) for \(\mathfrak A^*=\mathfrak B^*\),
and \(\mathfrak A\sim\mathfrak B\) for
\(\mathrm{Pow}(\mathfrak A)\cap\mathrm{Pow}(\mathfrak B)\neq\varnothing\).
By item~3 of \cref{thm:p03-w03-p03-018-keisler-compare}, which uses no
cardinal hypothesis, \(\mathfrak A\sim\mathfrak B\) iff
\(\mathfrak A\equiv\mathfrak B\).  Three facts are used repeatedly.
\begin{enumerate}
  \item[(a)] \(\mathfrak A\in\mathrm{Pow}(\mathfrak A)\): with
        \(I=\{0\}\) and the ultrafilter \(\{I\}\), which is principal,
        \(\mathfrak A^I/\{I\}\cong\mathfrak A\) by item~3 of
        \cref{thm:p03-w03-p03-018-comparison}.  Hence
        \(K\subseteq\mathrm{Pow}(K)\subseteq\mathrm{Prod}(K)\), the second
        inclusion because an ultrapower is the ultraproduct of a constant
        system.
  \item[(b)] Every member of \(\mathrm{Pow}(\mathfrak A)\) is elementarily
        equivalent to \(\mathfrak A\) (\L o\'s's theorem, \cref{thm:los},
        through \cref{cor:diagonal-ultrapower}, and
        \cref{lem:iso-invariance}); and an ultraproduct of a system of
        models of a sentence \(\sigma\) is a model of \(\sigma\), by
        \L o\'s's theorem, since the set of indices where \(\sigma\) holds
        is the whole index set.
  \item[(c)] (Keisler's Theorem~1.2.)  If \(K\in EC_{\Delta\Sigma}\), then
        \(K=\mathrm{Pow}(K)\): \(K\subseteq\mathrm{Pow}(K)\) by (a), and a
        member of \(\mathrm{Pow}(K)\) is elementarily equivalent to a member
        of \(K\) by (b), hence lies in \(K\), which is closed under
        elementary equivalence.
\end{enumerate}

\emph{Item~2.}  Let \(N\) be the elementarily closed class generated by
\(M\).  If \(\mathfrak A\in N\), then \(\mathfrak A\equiv\mathfrak B\) for
some \(\mathfrak B\in M\), so \(\mathfrak A\sim\mathfrak B\): some
structure lies in \(\mathrm{Pow}(\mathfrak A)\cap\mathrm{Pow}(\mathfrak B)
\subseteq\mathrm{Pow}(\mathfrak A)\cap\mathrm{Pow}(M)\).  Conversely, if
\(\mathfrak C\in\mathrm{Pow}(\mathfrak A)\cap\mathrm{Pow}(M)\), then
\(\mathfrak C\) is isomorphic to an ultrapower of some
\(\mathfrak B\in M\), and \(\mathfrak A\equiv\mathfrak C\equiv\mathfrak B\)
by (b), so \(\mathfrak A\in N\).

\emph{Item~1.}  (i)\(\Rightarrow\)(ii): \(K\) and \(\bar K\) are both
closed under elementary equivalence, so (c) applies to each.
(ii)\(\Rightarrow\)(iii): take \(M=K\).  If \(\mathfrak A\in K\), then
\(\mathfrak A\in\mathrm{Pow}(\mathfrak A)\cap\mathrm{Pow}(K)\) by (a).  If
\(\mathfrak C\in\mathrm{Pow}(\mathfrak A)\cap\mathrm{Pow}(K)\), then
\(\mathfrak C\in\mathrm{Pow}(K)=K\); were \(\mathfrak A\in\bar K\), also
\(\mathfrak C\in\mathrm{Pow}(\bar K)=\bar K\), which is impossible; so
\(\mathfrak A\in K\).  (iii)\(\Rightarrow\)(i): by item~2,
\(K=\{\mathfrak A:\mathrm{Pow}(\mathfrak A)\cap\mathrm{Pow}(M)\neq
\varnothing\}\) is the elementarily closed class generated by \(M\), which
is closed under elementary equivalence.

\emph{Item~3.}  Let \(N\) be the elementarily closed class generated by
\(K\); then \(K\subseteq N\).  If \(\mathfrak A\in M\cap N\), item~2 gives
\(\mathfrak C\in\mathrm{Pow}(\mathfrak A)\cap\mathrm{Pow}(K)\); then
\(\mathfrak C\in\mathrm{Pow}(M)=M\) and \(\mathfrak C\in\mathrm{Pow}(K)=K\),
contrary to \(K\cap M=\varnothing\).

\emph{Item~4.}  First a lemma, the content of Keisler's Theorem~1.3 that is
used here: \emph{if \(K\neq\varnothing\) and \(\mathfrak A\) is a model of
\(K^*\), then \(\mathfrak A\) is elementarily equivalent to an ultraproduct
of a system of members of \(K\).}  Let \(I\) be the set of finite subsets of
\(\mathfrak A^*\), and for \(s\in I\) let \(\sigma_s\) be the conjunction of
\(s\) (the sentence \(\forall v_0\,v_0=v_0\) when \(s=\varnothing\)).  Since
\(\mathfrak A\models\sigma_s\) and \(\mathfrak A\) is a model of \(K^*\),
\(\neg\sigma_s\notin K^*\), so some member of \(K\) satisfies \(\sigma_s\).
The members of \(K\) satisfying \(\sigma_s\) that have least rank in the
cumulative hierarchy form a nonempty set; by the Axiom of Choice pick one,
\(\mathfrak A_s\).  For \(s\in I\) put \(\hat s=\{t\in I:s\subseteq t\}\);
since \(\hat s\cap\hat t=\widehat{s\cup t}\neq\varnothing\), these sets have
the finite intersection property, and the Ultrafilter Lemma gives an
ultrafilter \(U\) on \(I\) containing all of them.  For
\(\varphi\in\mathfrak A^*\), \(\{t:\mathfrak A_t\models\varphi\}\supseteq
\widehat{\{\varphi\}}\in U\), so \(\prod_t\mathfrak A_t/U\models\varphi\) by
\L o\'s's theorem.  As \(\mathfrak A^*\) is complete, the ultraproduct and
\(\mathfrak A\) have the same theory.

Next, the direction of Keisler's Theorem~1.4 that is used here: \emph{if
\(K\) is closed under elementary equivalence, \(K=\mathrm{Prod}(K)\) and
\(\bar K=\mathrm{Prod}(\bar K)\), then \(K\in EC\).}  If \(K=\varnothing\),
\(K\) is the class of models of \(\neg\forall v_0\,v_0=v_0\); if
\(\bar K=\varnothing\), of \(\forall v_0\,v_0=v_0\).  Otherwise \(K\) is the
class of models of \(K^*\): a model of \(K^*\) is, by the lemma, elementarily
equivalent to a member of \(\mathrm{Prod}(K)=K\), hence in \(K\).  The class
\(\bar K\) is also closed under elementary equivalence, so in the same way it
is the class of models of \(\bar K^*\).  No structure is a model of
\(K^*\cup\bar K^*\), so by compactness
(\cref{thm:p03-compactness-ultraproduct}) there are finite
\(\Gamma_0\subseteq K^*\) and \(\Gamma_1\subseteq\bar K^*\) such that
\(\Gamma_0\cup\Gamma_1\) has no model.  Let \(\sigma\) be the conjunction of
\(\Gamma_0\) (or \(\forall v_0\,v_0=v_0\) if \(\Gamma_0\) is empty).  Every
member of \(K\) satisfies \(\sigma\).  A model \(\mathfrak B\) of \(\sigma\)
outside \(K\) would lie in \(\bar K\), hence satisfy \(\Gamma_1\) as well as
\(\Gamma_0\), which is impossible.  So \(K\) is the class of models of
\(\sigma\).

Now item~4.  If \(K\) is the class of models of \(\sigma\), then
\(K\subseteq\mathrm{Prod}(K)\) by (a), and \(\mathrm{Prod}(K)\subseteq K\)
by (b) together with isomorphism invariance; likewise for \(\bar K\), the
class of models of \(\neg\sigma\).  Conversely, suppose
\(K=\mathrm{Prod}(K)\) and \(\bar K=\mathrm{Prod}(\bar K)\).  By (a),
\(K\subseteq\mathrm{Pow}(K)\subseteq\mathrm{Prod}(K)=K\), so
\(K=\mathrm{Pow}(K)\), and likewise \(\bar K=\mathrm{Pow}(\bar K)\).  By
item~1, \(K\in EC_{\Delta\Sigma}\), that is \(K\) is closed under elementary
equivalence; by the previous paragraph \(K\in EC\).  No step uses a
hypothesis on cardinal exponentiation; the inputs are
\cref{thm:p03-w03-p03-017-keisler-shelah} (through item~3 of
\cref{thm:p03-w03-p03-018-keisler-compare}), \L o\'s's theorem, compactness
and the Axiom of Choice.  Keisler's own proof of Theorem~2.11 passes through
his Theorems~2.6 and~1.4 in the same order
\cite[p.~488]{Keisler1961ElementaryClasses}.
\end{proof}
\end{annexitem}
